\documentclass[11pt,a4paper]{article}
\usepackage[T1]{fontenc}
\usepackage[margin=1in]{geometry}
\usepackage{amsmath,amssymb,amsthm,mathtools}
\usepackage{booktabs}
\usepackage{xcolor}
\usepackage[colorlinks=true,linkcolor=blue!60!black,citecolor=blue!60!black,
            urlcolor=blue!60!black]{hyperref}
\usepackage{microtype}
\usepackage{enumitem}
\usepackage[most]{tcolorbox}

\newcommand{\lean}[1]{\texttt{\small #1}}
\newcommand{\dB}{d_{B}}
\newcommand{\Dgm}{\mathrm{Dgm}}
\newtcolorbox{keybox}[1][]{colback=blue!3,colframe=blue!45!black,
  boxrule=0.5pt,left=4pt,right=4pt,top=3pt,bottom=3pt,#1}
\newtcolorbox{failbox}[1][]{colback=red!3,colframe=red!45!black,
  boxrule=0.5pt,left=4pt,right=4pt,top=3pt,bottom=3pt,#1}
\newtcolorbox{fixbox}[1][]{colback=orange!4,colframe=orange!60!black,
  boxrule=0.5pt,left=4pt,right=4pt,top=3pt,bottom=3pt,#1}
\newtheorem*{theorem*}{Theorem}
\newcommand{\phaseprev}{0.6490}

\title{\textbf{How loose is Wasserstein stability on a vortex field?}\\
\large The Skraba--Turner cellular bound measured and certified on seven Gross--Pitaevskii pairs,\\
and where the duality points next}

\author{Xavier Callens\\
\small SocrateAI-Scientific-QuantumFluids\\
\small \url{https://github.com/xaviercallens/SocrateAI-Scientific-QuantumFluids}}
\date{September 2026}

\begin{document}
\maketitle

\begin{abstract}
\noindent
The previous report in this series~\cite{Callens2026Loop} found that the bottleneck stability theorem
for persistence diagrams, $\dB\le\|f-g\|_\infty$, holds but is empty on six of seven pairs of
two-dimensional Gross--Pitaevskii density fields: a few vortex-core pixels set the sup norm, and
$2\varepsilon$ exceeds the entire diagram. It named, as its first direction, the stability theorem in
which the sup norm is replaced by an $\ell_p$ norm, and set a gate: read the exact statement, fix every
convention in advance, compute the right-hand side before looking at the left. This report passes that
gate. The theorem is Skraba and Turner's cellular Wasserstein stability~\cite{SkrabaTurner2020}
(Theorem~4.6): $W_p(\Dgm f,\Dgm g)\le\|f-g\|_p$ with the norm taken over \emph{all} cells of the
complex, in the $\ell_p$ ground metric. We compute both sides exactly on the same seven pairs, for
$p=1,2$: the right-hand side as a finite sum over the cubical complex whose cell values are verified,
by an independent union-find, to be the ones GUDHI uses; the left-hand side as an optimal assignment on
the full diagrams, certified on every pair by explicit dual potentials (rigorous relative gap
$\le1.5\times10^{-4}$ on six certificates, $\le10^{-14}$ on the other twenty-two). Three controls pass exactly
before any pair is opened, one of them an exact known answer on real data (slack exactly $9$ for $p=1$
and $3$ for $p=2$, from the nine cells a raised minimum touches), and one of them caught an arithmetic
error in the author's own pre-registered hand value. The pre-registered question is not whether the
bound holds -- it is a theorem -- but whether it says anything the triangle inequality does not.
The answer is no, on every pair: the $\ell_1$ bound exceeds the trivial one by a factor $V_1$ of
$139$ (phase space) to $3200$ (the vortex-free control), the $\ell_2$ bound by $18$ to $65$, and the
measured slack $S_1=\|f-g\|_1/W_1$ runs from $490$ to $19{,}000$: in the theorem's own metric the
diagrams are three to four orders of magnitude closer than the fields. The pre-registered prediction
that the phase-space pair -- the one pair on which the bottleneck bound had been informative -- would
be informative here too, \emph{failed}. The reason is structural and is measured, not argued: the
$\ell_p$ norm is extensive in the number of cells, the diagram in the number of critical values, and
a resolution or time comparison changes every cell; the ratio of the two bounds scales as (cells per
bar)$^{1/p}$, which is why $V_2\approx\sqrt{V_1}$ on all seven pairs and why no $p$ helps -- small
$p$ pays the cell count, large $p$ returns to the vortex cores. The theorem is sharp exactly where the
control shows it: a perturbation that touches nine cells per critical value it moves. A second
finding was not looked for: the previous report's phase-space certificate had been computed on five
points per side, the same class of defect as the corruption it reported; it is corrected here. Two further directions are taken
to their gates and stopped there, with the reasons written down: the Berezinskii--Kosterlitz--Thouless
transition, which these fields carry but which needs a finite-temperature run this project does not
have; and the Fricke involution of an external review, which we now compare with Dolgachev's theorem
on the mirror moduli space of K3 surfaces read in full, and find to be the same involution without the
group -- a shadow, not an instance; the group theory behind that verdict (the Fricke matrix normalizes
$\Gamma_0(n)$; its axis action is this project's involution) is proved in Lean~4, as group theory.
\end{abstract}

\section{Introduction}

A stability theorem for persistence diagrams is a promise: if two functions are close, their diagrams
are close. The promise is only worth invoking if ``close'' on the right-hand side is a distance the data
actually make small. On vortex fields the sup norm is not such a distance: a vortex core is a hole of
depth $n_0$ a few pixels wide, and any comparison that moves a core by one pixel -- a change of
resolution, a step in time -- pays the full depth in $\|f-g\|_\infty$. The previous report measured this
on seven pairs: the bottleneck bound held on all seven with margins of $2$ to $16$, and excluded no bar
at all on six of them~\cite{Callens2026Loop}.

The same report ended with three directions on \emph{duality}, each behind a gate. The first was the
obvious one: use the theorem whose right-hand side is an $\ell_p$ norm, since an $\ell_p$ norm
averages over the field and a few core pixels no longer dominate it. The gate was to read the theorem
before computing, and that gate did work: the draft attributed the $\ell_p$ bound to Cohen-Steiner,
Edelsbrunner, Harer and Mileyko~\cite{CSEHM2010}, whose bound in fact keeps the sup norm on the right
(with a fractional exponent and a total-persistence constant), and the correction was made before
publication. The theorem that changes the norm is Skraba and Turner's~\cite{SkrabaTurner2020}, and it
is the one tested here.

The question this report pre-registers (\lean{docs/designs/WASSERSTEIN\_STABILITY\_PREREG.md}, committed
before any code) is therefore not whether $W_p\le\|f-g\|_p$ -- it is a theorem, and a violation would
be a bug in our pipeline -- but two measurable things about it on real vortex fields:
\begin{enumerate}[nosep]
\item \emph{Is it informative?} The triangle inequality alone gives
$W_p(\Dgm f,\Dgm g)\le W_p(\Dgm f,\emptyset)+W_p(\Dgm g,\emptyset)=:Z_p$, the cost of erasing both
diagrams. The theorem says something only where $\|f-g\|_p<Z_p$; we call $V_p=\|f-g\|_p/Z_p$ the
\emph{vacuity ratio}.
\item \emph{How loose is it?} The \emph{slack} $S_p=\|f-g\|_p/W_p$ measures how much closer the
diagrams are than the fields, in the theorem's own metric. It is a property of the data class, and
it is what one would need to know to turn the theorem into a resolution criterion.
\end{enumerate}

The second contribution is methodological and continues the previous report's own lesson. There, a
background agent corrupted a list of points between two pipeline stages, and only an independent
recomputation caught it. Here nothing moves between agents: one script computes, its controls read the
arrays from disk again, and the left-hand side is certified on every pair by a primal plan and dual
potentials checked by a function that knows nothing about how they were found.

\section{The theorem, with its conventions fixed}
\label{sec:theorem}

\begin{theorem*}[Skraba--Turner, Theorem 4.6, Cellular Wasserstein Stability~\cite{SkrabaTurner2020}]
Let $f,g:K\to\mathbb{R}$ be monotone functions on a finite CW complex $K$. Then
\begin{enumerate}[nosep,label=(\roman*)]
\item $W_p(\Dgm f,\Dgm g)\le\|f-g\|_p$, where $\|f-g\|_p^p=\sum_{\sigma\in K}|f(\sigma)-g(\sigma)|^p$;
\item $W_p(\Dgm_k f,\Dgm_k g)^p\le\sum_{\dim\sigma\in\{k,k+1\}}|f(\sigma)-g(\sigma)|^p$.
\end{enumerate}
\end{theorem*}

Three conventions travel with the statement and are adopted here as the theorem's own metric, not
ours. Distances in the plane are $\ell_p$; a point $(a,b)$ is at perpendicular distance
$2^{(1-p)/p}|b-a|$ from the diagonal ($|b-a|$ for $p=1$, $|b-a|/\sqrt2$ for $p=2$); essential classes
are matched among themselves at cost $|a-b|$ and never to a finite point or the diagonal; and the
total distance in (i) is $(\sum_k W_p(\Dgm_k f,\Dgm_k g)^p)^{1/p}$ over all degrees, each degree
matched separately. This is \emph{not} the $\ell^\infty$ ground cost of the previous report's
certificates; for $p=1$ the $\ell_1$ ground cost dominates it, so everything certified there is also
bounded by the theorem, but the comparison is made in the theorem's metric and the earlier numbers
appear only as a consistency check (prediction P7).

\paragraph{The complex and the cell values.} GUDHI's periodic cubical complex on a $2$D array with
\lean{top\_dimensional\_cells} is the $T$-construction~\cite{Bleile2021,GUDHI}: each pixel is a
$2$-cell with its own value and every lower cell takes the \emph{minimum} over the pixels containing
it. On an $N\times M$ torus that is $NM$ pixels, $2NM$ edges and $NM$ vertices; along a non-periodic
axis the boundary cells take the minimum over the pixels that exist. Both functions of every pair live
on one such complex (the resolution pairs by the restriction rule of the previous report: the fine
field restricted to the coarse grid), so the theorem applies literally, and $\|f-g\|_p$ is an exact
finite sum over $4NM$ numbers.

\section{Method and controls}
\label{sec:method}

\paragraph{Right-hand side.} \lean{cell\_values} builds the four cell families from the padded array
(wrap-around padding on a periodic axis, $+\infty$ padding otherwise) and \lean{cell\_norm} sums
$|f(\sigma)-g(\sigma)|^p$ over all of them, or over vertices and edges only for part~(ii) with $k=0$.

\paragraph{Left-hand side.} For each degree, the augmented $(n+m)\times(m+n)$ matrix of
Kerber--Morozov--Nigmetov~\cite{Kerber2017} with $p$-th powers of $\ell_p$ costs (each point owning one
diagonal slot; cross-use forbidden by a large penalty; dummy-to-dummy free) is solved as an assignment
(\lean{scipy.optimize.linear\_sum\_assignment}); the sum is the $p$-th power of $W_p^{(k)}$. Then the
dual LP $\max\sum u+\sum v$ s.t.\ $u_i+v_j\le C_{ij}$ is solved (HiGHS, sparse; the dummy block
replaced by two auxiliary variables and the penalty constraints dropped, which \emph{relaxes} the dual
and is therefore not trusted) and an independent function checks the potentials against the
\emph{full} matrix: feasibility everywhere and tightness against the plan. A plan with feasible, tight
potentials is optimal by finite LP weak duality -- the theorem proved in general in
\lean{lean\_src/WassersteinCertificate.lean} of the previous report, which covers every certificate
here. The diagrams are \emph{full}: every finite bar and every essential class, not the top-$k$
truncation used before.

\paragraph{Controls, all passed before any pair was opened.}
\begin{itemize}[nosep]
\item \textbf{C1} (the cell values are GUDHI's). An independent union-find over the vertices and edges
we build, with the elder rule, must reproduce GUDHI's $\Dgm_0$ exactly. Tested on random small arrays
under all four periodicity patterns, with and without ties, and then on all twelve arrays of the seven
pairs. All exact.
\item \textbf{C2} (hand-computed known answer). The eight-point cycle of the previous report,
$f=(1,5,2,6,3,7,4,8)$ and $g$ with $v_6:=4.5$. Only one pixel changes and neither neighbouring vertex
does (their minima are $3$ and $4$ on both sides), so $\|f-g\|_1=\|f-g\|_2=5/2$; the diagrams differ in
one point, $(4,7)\to(4,4.5)$, so $W_1=5/2$ and the $p=1$ bound is \emph{tight}. The script also had to
return the previous report's $\ell^\infty$-ground $W_1=7/4$ on the same example when told to use that
metric, as the check that the two conventions are the ones we think they are. It did.
\begin{fixbox}
\textbf{Amendment A1, dated before any pair.} The pre-registered hand value for $W_2$ on this example
was $5/2$. The first control run returned $2.1506\ldots$, and a second hand computation agrees: for
$p=2$ sending both points to the diagonal costs $\big((3/\sqrt2)^2+(1/2\sqrt2)^2\big)^{1/2}
=\sqrt{37/8}$, which is cheaper than the $5/2$ of matching them. The expected value was corrected to
$\sqrt{37/8}$ and nothing else changed. The $p=2$ bound holds on this example and is not tight. The
control did what a control is for: it caught the author, not the code.
\end{fixbox}
\item \textbf{C3} (exact known answer on real data). Take $\rho$ of pair T1's first field, the birth
pixel of its deepest finite $H_0$ bar that is a strict local minimum in $8$-connectivity, and raise it
by $\delta$. All nine cells that pixel touches (itself, four edges, four vertices) have it as their
minimum and change by exactly $\delta$: $\|f-g\|_1=9\delta$, $\|f-g\|_2=3\delta$, the part-(ii) $k=0$
sums are $8\delta$ and $\sqrt8\,\delta$, and the diagram moves one birth by $\delta$: $W_1=W_2=\delta$.
Measured: slack exactly $9.000$ and $3.000$. The pre-registration allowed $\delta$ to be reduced if it
exceeded the gap to the pixel's nearest neighbour; it did ($\delta=0.01\times$depth exceeded the gap
$4.3\times10^{-3}$) and $\delta$ was set to half the gap, as written, and recorded.
\item \textbf{C4} (negative). Pair T1's optimal plan with pair T2's potentials must be rejected by the
certificate checker. It was: maximal violation $0.898$, primal $10.56$ against dual $21.60$, not
feasible, not tight.
\end{itemize}

\section{Results}
\label{sec:results}

Table~\ref{tab:main} gives, for each pair and $p\in\{1,2\}$: the theorem's right-hand side
$\|f-g\|_p$ over all cells; the part-(ii) bound for $k=0$; the total $W_p$ and its $H_0$ part
$W_p^{(0)}$; the trivial bound $Z_p$; the slack $S_p$ and the vacuity ratio $V_p$.

\begin{table}[h]\centering\small
\begin{tabular}{llrrrrrrrc}
\toprule
pair & $p$ & $\|f-g\|_p$ & $B_p^{(ii,0)}$ & $W_p$ & $W_p^{(0)}$ & $Z_p$ & $S_p$ & $V_p$ & bound \\
\midrule
R1 & 1 & 2.613e+05 & 1.964e+05 & 45.67 & 19.56 & 698.9 & 5.72e+03 & 374 & holds \\
R1 & 2 & 364.9 & 316.8 & 1.47 & 0.979 & 20.13 & 248 & 18.1 & holds \\
R2 & 1 & 3.712e+05 & 2.788e+05 & 55.46 & 23 & 868.7 & 6.69e+03 & 427 & holds \\
R2 & 2 & 476.4 & 413.3 & 1.513 & 0.8816 & 21.95 & 315 & 21.7 & holds \\
R3 & 1 & 4.847e+05 & 3.633e+05 & 79.58 & 28.79 & 1291 & 6.09e+03 & 376 & holds \\
R3 & 2 & 593.1 & 513.3 & 1.716 & 0.8722 & 26.62 & 346 & 22.3 & holds \\
T1 & 1 & 1.99e+06 & 1.495e+06 & 103 & 36.83 & 794.9 & 1.93e+04 & 2.5e+03 & holds \\
T1 & 2 & 1205 & 1045 & 2.666 & 1.459 & 21.3 & 452 & 56.6 & holds \\
T2 & 1 & 2.102e+06 & 1.578e+06 & 248 & 84.23 & 1109 & 8.48e+03 & 1.9e+03 & holds \\
T2 & 2 & 1276 & 1106 & 4.722 & 2.343 & 24.82 & 270 & 51.4 & holds \\
S1 & 1 & 1.695e+06 & 1.274e+06 & 285.5 & 167.5 & 530.4 & 5.94e+03 & 3.2e+03 & holds \\
S1 & 2 & 1015 & 880.5 & 8.117 & 6.343 & 15.63 & 125 & 64.9 & holds \\
P1 & 1 & 9707 & 7179 & 19.7 & 6.445 & 70.06 & 493 & 139 & holds \\
P1 & 2 & 30.31 & 25.88 & 0.3871 & 0.2349 & 1.645 & 78.3 & 18.4 & holds \\
\bottomrule
\end{tabular}

\caption{The Skraba--Turner bound on the seven pairs, in its own metric ($\ell_p$ ground cost, full
diagrams, all degrees). $S_p=\|f-g\|_p/W_p$; $V_p=\|f-g\|_p/Z_p$, informative iff $V_p<1$.}
\label{tab:main}
\end{table}

\begin{table}[h]\centering\small
\begin{tabular}{llrrrlr}
\toprule
pair & $p$ & rel.\ gap & max viol.\ $\delta$ & certified gap & scope & LP s \\
\midrule
R1 & 1 & 0.0e+00 & 1.1e-16 & 8.5e-15 & full & 29.5 \\
R1 & 2 & 1.8e-07 & 9.6e-08 & 1.5e-04 & full & 384.0 \\
R2 & 1 & 0.0e+00 & 7.8e-17 & 6.2e-15 & full & 65.1 \\
R2 & 2 & 3.3e-16 & 3.5e-18 & 7.9e-15 & full & 62.2 \\
R3 & 1 & 1.3e-16 & 1.4e-17 & 9.0e-16 & top-200 & 1.2 \\
R3 & 2 & 1.1e-15 & 1.7e-18 & 1.1e-15 & top-200 & 1.3 \\
T1 & 1 & 1.9e-16 & 9.7e-17 & 4.4e-15 & full & 60.6 \\
T1 & 2 & 2.2e-07 & 1.4e-07 & 1.2e-04 & full & 834.3 \\
T2 & 1 & 4.9e-16 & 3.3e-16 & 6.7e-15 & top-200 & 2.3 \\
T2 & 2 & 0.0e+00 & 1.4e-17 & 2.2e-15 & top-200 & 1.4 \\
S1 & 1 & 1.7e-16 & 3.3e-16 & 2.7e-15 & full & 27.4 \\
S1 & 2 & 0.0e+00 & 2.8e-17 & 8.8e-16 & full & 24.3 \\
P1 & 1 & 2.6e-08 & 3.9e-08 & 4.2e-06 & top-200 & 3.9 \\
P1 & 2 & 4.5e-09 & 6.8e-10 & 5.0e-06 & top-200 & 3.6 \\
\bottomrule
\end{tabular}

\caption{Duality certificates for $W_p^{(0)}$: relative primal--dual gap as returned, maximal
violation $\delta$ of $u_i+v_j\le C_{ij}$ over the full augmented matrix, the rigorous certified
relative gap after the $-R\delta$ shift, the scope (full diagram, or the 200 longest bars per side
where the full LP exceeded the size rule of amendment A2), and LP time.}
\label{tab:cert}
\end{table}

\paragraph{Pre-registered verdicts.}
\begin{itemize}[nosep]
\item \textbf{P1} (the bound holds, parts (i) and (ii)): holds on all seven pairs for $p=1,2$, by
factors of $S_p$. As pre-registered, this line tests the pipeline, not the theorem.
\item \textbf{P2} (resolution pairs vacuous at $p=1$): \textbf{PASS} -- $V_1=374,\ 427,\ 376$.
\item \textbf{P3} (resolution pairs vacuous at $p=2$ as well): \textbf{PASS} -- $V_2=18.1,\ 21.7,\
22.3$. With P2, the kill of the previous report's \S7.1 applies: no Wasserstein stability theorem
makes the $\xi/\Delta x=4$ vs.\ $8$ comparison informative.
\item \textbf{P4} (phase-space pair informative at $p=1$): \textbf{FAIL} -- $V_1=139$, $V_2=18.4$.
The pair on which the bottleneck bound was informative ($2\varepsilon=0.58$ against a persistence
range of about $0.2$ was the previous report's finding; here $\varepsilon=0.2878$ is reproduced) is
not informative under the $\ell_p$ bound. See the mechanism below.
\item \textbf{P5} (time pairs and the control, reported): $V_1=2500,\ 1900,\ 3200$ for T1, T2, S1;
these are the largest, because the fields differ everywhere.
\item \textbf{P6} (slack $S_1\ge10$ on every GP pair): \textbf{PASS}, by a wide margin --
$S_1$ from $5720$ (R1) to $19{,}300$ (T1); $S_2$ from $125$ to $452$.
\item \textbf{P7} (consistency with the previous report's $\ell^\infty$ top-30 $W_1$): \textbf{PASS}
on all seven; the six GP values are reproduced to all printed digits ($2.3371$, $0.5791$, $0.8031$,
$2.4532$, $0.9513$, $10.8003$), and $\varepsilon$ likewise. The phase-space value is discussed in
\S\ref{sec:erratum}.
\end{itemize}

\paragraph{Why every pair is vacuous, measured.} Write $\bar\Delta$ for the mean of
$|f(\sigma)-g(\sigma)|$ over the $4NM$ cells and $\bar\pi$ for the mean persistence of the
$n+m$ finite bars. Then $\|f-g\|_1=4NM\,\bar\Delta$ and $Z_1=(n+m)\,\bar\pi$, so
$V_1=\big(4NM/(n+m)\big)\cdot\big(\bar\Delta/\bar\pi\big)$: the number of cells per bar times
the ratio of mean cell change to mean bar length. On R1, $4NM/(n+m)=1.05\times10^6/1495=701$ and the
measured $\bar\Delta=0.249\,n_0$ -- the two resolutions are two different turbulent trajectories,
not one field sampled twice, and they differ by a quarter of the mean density on the \emph{average}
cell. On T1 the mean cell change is $0.47\,n_0$; on the vortex-free control S1, $0.40\,n_0$. Nothing
about vortex cores enters: the $\ell_1$ bound is empty because the perturbation is dense, and a
dense perturbation of a field with $10^3$ cells per feature is charged $10^3$ times per feature.

For general $p$ the same count gives $V_p\sim\big(4NM/(n+m)\big)^{1/p}\,\bar\Delta_p/\bar\pi_p$,
and the data follow it: $V_2/\sqrt{V_1}=0.94,\ 1.05,\ 1.15,\ 1.13,\ 1.18,\ 1.15,\ 1.56$ across the
seven pairs. The limit $p\to\infty$ is the bottleneck case of the previous report, vacuous for the
other reason (a few core pixels set the norm, and $2\varepsilon$ exceeds the diagram). There is no
$p$ at which a dense perturbation of a vortex field is small in $\ell_p$ and large in persistence:
small $p$ pays the cell count, large $p$ pays the core depth.

\paragraph{Exploratory: is there a $p$ that helps?} Not pre-registered, computed after the verdicts
above and labelled accordingly: the $H_0$ vacuity ratio $V_p=\|f-g\|_p/Z^{(0)}_p$ (all cells over the
$H_0$ diagrams alone, primal only, no certificate) for $p\in\{1,2,4,8\}$ and, at $p=\infty$, the
bottleneck case $\varepsilon/(\max\text{persistence}/2)$ of the previous report.

\begin{center}\small
\begin{tabular}{lrrrrr}
\toprule
pair & $V_1$ & $V_2$ & $V_4$ & $V_8$ & $V_\infty$ \\
\midrule
R1 & 694 & 24.2 & 5.1 & 2.62 & 2.95 \\
R2 & 869 & 31.1 & 6.12 & 2.94 & 3.52 \\
R3 & 855 & 35.1 & 6.96 & 3.21 & 3.67 \\
T1 & 4.92e+03 & 78.8 & 9.84 & 3.6 & 3.02 \\
T2 & 4.19e+03 & 79.1 & 10.6 & 4.02 & 3.71 \\
S1 & 6.1e+03 & 87.2 & 10.2 & 3.49 & 2.06 \\
P1 & 387 & 30.5 & 5.95 & 2.4 & 3 \\
\bottomrule
\end{tabular}

\end{center}

$V_p$ falls with $p$ as the cell count is taken to the power $1/p$, reaches $2.4$--$4.0$ at $p=8$,
and is $2.1$--$3.7$ at $p=\infty$. Its minimum over $p$ is above $2$ on every pair. The two vacuity
mechanisms -- the cell count at small $p$, the core depth at large $p$ -- meet in the middle without
ever crossing $1$. This is the measured form of the sentence in the abstract: on a dense perturbation
of a vortex field, no $\ell_p$ stability bound says anything the triangle inequality does not.

\paragraph{What the slack measures.} $S_1$ is between $5\times10^3$ and $2\times10^4$ on the GP
pairs and $493$ on the phase-space pair. The diagrams are not far apart: $W_1/Z_1$, the actual
distance as a fraction of the cost of erasing both diagrams, is $0.065$ (R1), $0.064$ (R2), $0.062$
(R3), $0.13$ (T1), $0.22$ (T2), $0.54$ (S1), $0.28$ (P1). The resolution pairs' diagrams are within
$7\%$ of erasure of each other -- close, by the standard the theorem uses -- while their fields are
not close at all. That is a statement about persistence as a summary, and it is the quantitative
content of this round: on this data class the $H_0$ diagram is stable to dense perturbations to a
degree the stability theorem cannot see, because the theorem is sharp only for sparse ones. The
control C3 shows the sharp case exactly: slack $9$ at $p=1$ and $3$ at $p=2$, the coordination number
of a pixel and its square root.

\paragraph{Part (ii) against part (i).} Restricting to vertices and edges removes the pixel term
and nothing else; on every pair $B^{(ii,0)}_1/B^{(i)}_1$ lies between $0.74$ and $0.75$ ($3/4$ would
be exact if the three cell families changed equally), and $B^{(ii,0)}_2/B^{(i)}_2$ between $0.85$ and $0.87$. The
per-degree refinement is real and small; it does not change any verdict.

\section{Where the duality points next}
\label{sec:directions}

\subsection{What this bound can and cannot become}

The kill written into the previous report is now executed: the resolution comparison
$\xi/\Delta x=4$ vs.\ $8$ is not a perturbation in any $\ell_p$ sense, and no stability theorem of
this family rescues it. The honest use of persistence at $\xi/\Delta x=4$ is therefore not ``the
diagram is guaranteed close to the converged one'' but ``the diagram is measured to be within $7\%$
of erasure of the converged one, with no guarantee''. Where the theorem \emph{is} sharp -- sparse
perturbations, C3 -- it is a tool for something else: certifying that a localized edit to a field
(a planted vortex moved, a core filled in) changes the diagram by at most nine cells' worth, which is
a statement one could use in a control experiment. That is the direction the bound can become, and it
is small. The number to carry forward is the slack: on vortex fields, $10^3$--$10^4$ at $p=1$.

\subsection{Erratum to the previous report}
\label{sec:erratum}

The previous report's Table~3 lists the phase-space pair's certificate as $n,m=5,5$ with
$W_1=0.0632\ldots$. The full $H_0$ diagrams of the raw phase-space arrays have $2150$ and $1623$
finite bars, so the pre-registered top-30 truncation should have had 30 points per side. The
diagrams stage of that round was correct -- the bottleneck distance recomputed here on the full
diagrams is $0.0325030\ldots$, the published $\dB=0.0325$, and $\varepsilon=0.2878$ is reproduced --
and the $5\times5$ certificate is valid for the five points it was given; but the number is not the
$\ell^\infty$ top-30 $W_1$ of the phase-space diagrams. It is the same class of defect as the S1
corruption that report found and fixed, on the one pair it did not independently recompute. The
correct value, computed on the full diagrams in this round with the previous convention, is
$W_1^{\ell^\infty,\,\mathrm{top}\text{-}30}=\phaseprev$. The published PDF is not changed; the
erratum is recorded in the project ledger and here.

\subsection{The transition these fields carry: BKT, and the run this project does not have}

A vortex field in two dimensions is the setting of the Berezinskii--Kosterlitz--Thouless transition,
where bound vortex--antivortex pairs unbind and the superfluid density jumps by the universal amount
of Nelson and Kosterlitz~\cite{NelsonKosterlitz1977}, measured in helium films by Bishop and
Reppy~\cite{BishopReppy1978} and in a trapped atomic gas by Hadzibabic et al.~\cite{Hadzibabic2006}.
Persistence of the density field sees vortex cores as $H_0$ bars and, by the duality of
constructions~\cite{Bleile2021} used as a control in~\cite{Callens2026TDA}, nothing else; whether the
\emph{pairing} of cores -- the quantity that changes at the transition -- leaves a signature in
$\Dgm_1$ of the density or in the diagram of the phase field is a question with a known answer to test
against: the critical density $n_c\lambda_T^2=\ln(\xi/\tilde g)$ with $\xi\approx380$ of Prokof'ev,
Ruebenacker and Svistunov~\cite{Prokofev2001}.

\emph{Gate, and why it is not passed here.} The fields of this and the previous report are
zero-temperature Gross--Pitaevskii evolutions of planted vortices. BKT is a thermal transition; the
standard numerical route is the projected Gross--Pitaevskii (classical-field) method, used for exactly
this question by Simula and Blakie~\cite{SimulaBlakie2006} and by Foster, Blakie and
Davis~\cite{FosterBlakieDavis2010}, who resolve vortex-pair unbinding, the superfluid fraction and the
first-order correlation function across the transition in a homogeneous box. This project has no
finite-temperature solver. Writing one is a project of its own, with its own known answers (the
superfluid jump, $n_c$), and it must be pre-registered as such; nothing in the present data can be
reinterpreted as thermal. The direction stays open and unstarted.

\subsection{The Fricke involution, read against Dolgachev}
\label{sec:fricke}

The external review quoted in~\cite{Callens2026Loop} placed this project's older involution
$R\leftrightarrow\alpha'/R$ -- in \lean{lean\_src/DualLength.lean}, the map $k\mapsto k_s^2/k$ on
wavenumbers, which exchanges the phonon and free-particle terms of the Bogoliubov dual length
$\ell(k)=1/k+k/k_s^2$ and fixes $k=k_s$ -- with the Fricke involution and the Dolgachev--Nikulin mirror
symmetry of K3 surfaces, as ``a theorem'' that ``keeps the duality''. The gate set there was: read
Dolgachev and decide whether the inversion is an instance of his theorem or its shadow. We read
Dolgachev's paper~\cite{Dolgachev1996} in full (arXiv:alg-geom/9502005). The relevant statements are:

\begin{keybox}
\textbf{Theorem 7.1.} For the mirror family of degree-$2n$ polarized K3 surfaces (the $M_n$-polarized
surfaces, $M_n=U\perp\langle -2n\rangle$-orthogonal, Picard rank 19), the moduli space is
$K_{M_n}\cong\mathbb{H}/\Gamma_0(n)^+$, where $\Gamma_0(n)^+$ is $\Gamma_0(n)$ extended by the Fricke
involution $F=\begin{psmallmatrix}0&-1/\sqrt n\\ \sqrt n&0\end{psmallmatrix}$, $t\mapsto -1/(nt)$.
\textbf{Theorem 7.6.} An $M_n$-polarized K3 surface with period $t$ is, up to a canonical involution,
the Kummer surface of $E_t\times E_{-1/nt}$, a pair of $n$-isogenous elliptic curves; the Fricke
involution acts on $\mathbb{H}/\Gamma_0(n)$ by $(E,A)\mapsto(E/A,E_n/A)$.
\end{keybox}

Now compare. Restrict $F$ to the imaginary axis, $t=iy$: $F(iy)=i/(ny)$, i.e.\ $y\mapsto 1/(ny)$, an
inversion of the positive reals about $y=1/\sqrt n$. With $y=k$ and $n=1/k_s^2$ this is exactly
\lean{DualLength}'s $k\mapsto k_s^2/k$ with fixed point $k_s$. So the reviewer is right that it is the
same involution -- and so is every inversion $x\mapsto c/x$, which is the point. What makes Dolgachev's
statement a theorem about K3 surfaces is not $F$ but the group $\Gamma_0(n)$ that $F$ normalizes: the
translation $t\mapsto t+1$, the congruence condition $n\mid c$, the lattice-polarized family whose
periods live in $\mathbb{H}$, and the isogeny that $F$ realizes on the elliptic curves. A wavenumber $k$
carries none of it: there is no $t\mapsto t+1$ on the Bogoliubov dispersion, no lattice, no family
whose monodromy is $\Gamma_0(n)$. \lean{DualLength.lean} proves, and says it proves, an AM--GM
statement about two positive terms.

\emph{Verdict at the gate: shadow, not instance.} The honest sentence is ``the same involution, without
the group''. It is not a claim about helium, and no Lean line is written for it. What \emph{would} be
an instance is a system whose parameter space carries the modular action -- the quantum Hall
$\sigma$-plane under $\Gamma_0(2)$~\cite{LutkenRoss1992}, the Seiberg--Witten $u$-plane under
$\Gamma(2)$~\cite{SW1994} -- and those are the reviewer's other examples, outside this project's data.
Formalizing Dolgachev's Theorem~7.1 itself in Lean~4 is out of reach of Mathlib as pinned here (no
lattice-polarized K3 surfaces, no period map). The group-theoretic core is not, and it is now in
\lean{lean\_src/Fricke.lean} (12 theorems, axiom footprint $\{$\lean{propext},
\lean{Classical.choice}, \lean{Quot.sound}$\}$, three deliberately broken variants fail to compile),
labelled as group theory and nothing else: the integer Fricke matrix $W=\begin{psmallmatrix}0&-1\\
n&0\end{psmallmatrix}$ squares to $-n\cdot1$ (\lean{W\_mul\_W}); for every $A$ in Mathlib's
\lean{Gamma0 n} there is an explicit $B\in$ \lean{Gamma0 n} with $WA=BW$ (\lean{fricke\_normalizes}),
and $A\mapsto B$ is itself an involution (\lean{conj\_conj}) -- which is what makes $\Gamma_0(n)^+$ a
group containing $\Gamma_0(n)$ with index~2; the Möbius action of $W$ on the imaginary axis is
$iy\mapsto i/(ny)$ (\lean{fricke\_on\_axis}); with $n=1/k_s^2$ that map is literally
\lean{DualLength}'s $k\mapsto k_s^2/k$ (\lean{axis\_eq\_dual}), leaves $\ell(k)$ invariant
(\lean{ell\_axis\_invariant}) and has the unique positive fixed point $1/\sqrt n=k_s$
(\lean{axis\_fixed\_iff}). That is the whole content of ``the same involution'': the file proves what
$W$ does to $\Gamma_0(n)$ and what its shadow does to a wavenumber, and contains no statement about
K3 surfaces, mirror symmetry or helium, because there is none to prove.

\subsection{The certificate duality itself}

The one duality that has done work in two reports is the humblest: finite LP weak duality, which
turns a solver's answer into a proof. Its next step is not new physics but a larger reach: the seven
full-diagram certificates here are checked in Python by an independent function; the general theorem
is in Lean; the specific $1600\times1600$ floating-point matrices are not re-verified by the kernel.
Exact rational certificates for full diagrams, checked by \lean{norm\_num} or by a verified checker,
would close that last gap and are the natural P6 of a future round.

\section{Discussion}

Three directions were named behind gates; this report is the record of taking each to its gate.
The first passed and was measured: the cellular Wasserstein bound holds, as it must, and is empty on
every pair, with the mechanism quantified (cells per bar) and a control showing where it is sharp.
The pre-registered prediction that the phase-space pair would be informative was wrong, and is
reported as wrong. The second and third stopped at their gates -- one for want of a finite-temperature
solver, one because the theorem it was compared with turned out to be about a group the data does not
carry -- and the reasons are written down, which is what a gate is for.

What this round did not do. One seed, one solver, one data class; $p=1$ and $2$ only, with the
$p$-scaling inferred from two points and a count, not from a sweep; certificates on the full diagrams
where the LP fit in memory and on the 200 longest bars where it did not, scope recorded per line of
Table~\ref{tab:cert}; no Lean line for Theorem~4.6 itself, whose statement is about persistence
modules on a CW complex and is beyond the pinned library (the Lean written this round,
\lean{Fricke.lean}, is the group theory of \S\ref{sec:fricke}, not the topology of \S\ref{sec:theorem}). And the round's own small lesson: the
author's hand value for $W_2$ on an eight-point example was wrong, and the control caught it before
any data was opened -- the same discipline that found the previous report's five-point certificate
after publication rather than before. Redoing independently remains the only check that has worked.

\paragraph{Reproducibility.} \lean{exploration/tda/wp\_stability.py} runs the controls and the seven
pairs and writes \lean{data/generated/kinetic\_tda/wp\_stability\_results.json};
\lean{wp\_stability\_report.py} produces the tables above from that file and re-derives the verdicts.
Pre-registration and amendment: \lean{docs/designs/WASSERSTEIN\_STABILITY\_PREREG.md}. The data arrays
are the previous report's (\lean{make\_gp\_frames.py}, \lean{make\_gp\_frames\_n512.py}, seed as
recorded there).

\bibliographystyle{plain}
\bibliography{refs_wasserstein_slack}
\end{document}
